
@@M0@@
Comme @@M1@@ est inversible, on a @@M2@@, donc @@M3@@.
Par suite, on a :
@@M4@@

Par ailleurs, la fonctorialité de @@M5@@, et la formule de projection 3.7 entraînent :
@@M6@@
Or @@M7@@. Donc :
@@M8@@

Comme @@M9@@, @@M10@@ est l'image dans @@M11@@ de @@M12@@, donc c'est @@M13@@. Par ailleurs, on a :
@@M14@@

On obtient donc finalement :
@@M15@@
ce qui prouve la relation à démontrer.

\section*{5. Théorème de Riemann--Roch : cas du morphisme structural d'un fibré projectif}

\textbf{5.1.} D'après 3.9, il reste, pour prouver le théorème de Riemann--Roch, à le prouver pour le morphisme structural @@M16@@, où @@M17@@ est un schéma quasi-compact, admettant un faisceau ample. On pose @@M18@@, @@M19@@, et @@M20@@.

